\documentclass[11pt,a4paper]{article}

% Packages standards
\usepackage[utf8]{inputenc}
\usepackage[T1]{fontenc}
\usepackage[french]{babel}
\usepackage{amsmath, amssymb, amsfonts}
\usepackage{geometry}
\geometry{hmargin=2.5cm,vmargin=3cm}
\usepackage{booktabs} % Pour des tableaux professionnels
\usepackage{tikz} % Pour les diagrammes et graphes
\usetikzlibrary{arrows.meta, positioning, shapes.geometric}
\usepackage{hyperref}
\hypersetup{
    colorlinks=true,
    linkcolor=blue,
    filecolor=magenta,      
    urlcolor=cyan,
    pdftitle={Rapport d'Aide à la Décision Multicritère - Rock en Seine},
    pdfpagemode=FullScreen,
}
\usepackage{float}
\usepackage{enumitem}

\begin{document}

\begin{titlepage}
    \centering
    \vspace*{1.5cm}
    
    {\large\scshape Université Paris-Dauphine - PSL}\\[0.5cm]
    {\large\scshape Cours d'Aide Multicritère à la Décision}\\[1.5cm]
    
    \rule{\linewidth}{0.8mm}\\[0.6cm]
    {\huge\bfseries PROJET D'AIDE À LA DÉCISION COLLECTIVE}\\[0.4cm]
    {\large\itshape Application comparée de méthodes d'agrégation et de surclassement : \\ Le cas du choix des billets pour Rock en Seine}\\[0.6cm]
    \rule{\linewidth}{0.8mm}\\[2.5cm]
    
    \begin{minipage}{0.45\textwidth}
        \begin{flushleft} \large
            \emph{Auteurs :}\\
            Jules \textsc{Bobo}\\
            Florent \textsc{Jin}\\
            Victor \textsc{Cai}
        \end{flushleft}
    \end{minipage}
    ~
    \begin{minipage}{0.45\textwidth}
        \begin{flushright} \large
            \emph{Enseignante :} \\
            Lucie \textsc{Galand}\\
            \vspace{0.4cm}
            \emph{Date de rendu :} \\
            10 juin 2026
        \end{flushright}
    \end{minipage}
    
    \vfill
    
    \begin{center}
        \rule{0.5\linewidth}{0.3mm}\\[0.4cm]
        \begin{minipage}{0.85\linewidth}
            \centering
            \small\itshape
            \textbf{Résumé :} Ce rapport présente l'application d'un processus d'aide multicritère à la décision collective (AMDC) visant à résoudre un problème concret : le choix des jours de participation au festival Rock en Seine pour un groupe de trois amis. Nous modélisons les préférences divergentes (Jules préférant le vendredi ou le dimanche, Florent le mercredi et Victor le jeudi) sous forme de critères. Nous appliquons et comparons deux approches théoriques fondamentales étudiées en cours : la méthode d'agrégation de la Somme Pondérée (procédure compensatoire) et la méthode de surclassement ELECTRE I (procédure non-compensatoire). Les résultats mettent en évidence l'existence de trois solutions stables distinctes (Mercredi, Jeudi et Samedi) sous forme de noyaux multiples, permettant ainsi d'éclairer la négociation collective du groupe.
        \end{minipage}
    \end{center}
    
\end{titlepage}

\newpage
\tableofcontents
\newpage

% ========================================== %
% SECTION 1: INTRODUCTION ET MODÉLISATION   %
% ========================================== %
\section{Présentation et modélisation du problème}
\label{sec:presentation}

\subsection{Contexte de la décision collective}
Le choix de participer à un festival de musique comme Rock en Seine entre amis est une situation typique de décision collective multicritère. Le groupe est composé de trois décideurs : Jules ($D_1$), Florent ($D_2$) et Victor ($D_3$). Bien que le désir d'assister au festival soit partagé, des désaccords majeurs apparaissent concernant les jours de présence. Ces désaccords reposent sur plusieurs dimensions :
\begin{itemize}
    \item \textbf{La programmation artistique} : Chaque journée du festival propose une liste d'artistes différente. Les goûts musicaux des trois décideurs divergent : Jules préfère la programmation du vendredi ou du dimanche, Florent adore celle du mercredi et Victor a une nette préférence pour celle du jeudi.
    \item \textbf{Le coût financier} : Le prix unitaire d'un billet d'une journée est de 75~€, tandis que le Pass 3 jours est proposé au tarif de 185~€ (économie de 40~€).
    \item \textbf{Les disponibilités temporelles} : Des contraintes professionnelles ou personnelles limitent la liberté de chacun (Florent travaille le vendredi soir, Victor a un repas de famille obligatoire le dimanche).
\end{itemize}

L'objectif de cette étude est d'appliquer un processus d'aide à la décision afin de recommander la formule de billet la plus adaptée au groupe (problématique de choix, notée $P.\alpha$).

\subsection{Identification des alternatives (Ensemble des actions $A$)}
Nous identifions les alternatives comme étant les différentes formules de billets ou combinaisons de jours envisageables pour le groupe. Compte tenu des préférences exprimées, nous considérons l'ensemble des actions $A = \{a_1, a_2, a_3, a_4, a_5, a_6\}$ défini de la manière suivante :
\begin{itemize}
    \item $a_1$ : Billet 1 Jour - Mercredi (Me)
    \item $a_2$ : Billet 1 Jour - Jeudi (Je)
    \item $a_3$ : Billet 1 Jour - Vendredi (Ve)
    \item $a_4$ : Billet 1 Jour - Samedi (Sa) (journée intermédiaire de consensus)
    \item $a_5$ : Billet 1 Jour - Dimanche (Di)
    \item $a_6$ : Pass 3 Jours - Vendredi + Samedi + Dimanche (3 Jours)
\end{itemize}

\subsection{Famille de critères d'évaluation ($F$)}
Pour évaluer et comparer ces alternatives de manière structurée et lisible, nous définissons une famille cohérente de $p = 7$ critères, notée $F = \{g_1, g_2, g_3, g_4, g_5, g_6, g_7\}$ :
\begin{enumerate}
    \item \textbf{Coût financier ($g_1$)} : Prix unitaire du billet ou pass en euros (critère quantitatif à minimiser).
    \item \textbf{Satisfaction artistique Jules ($g_2$)} : Appréciation qualitative de la programmation musicale par Jules sur une échelle de 0 à 10, à maximiser.
    \item \textbf{Satisfaction artistique Florent ($g_3$)} : Appréciation de la programmation par Florent (échelle 0-10, à maximiser).
    \item \textbf{Satisfaction artistique Victor ($g_4$)} : Appréciation de la programmation par Victor (échelle 0-10, à maximiser).
    \item \textbf{Disponibilité Jules ($g_5$)} : Facilité logistique et temporelle pour Jules d'assister au festival pour la formule choisie, évaluée sur une échelle de 0 (impossible) à 5 (parfaite disponibilité), à maximiser.
    \item \textbf{Disponibilité Florent ($g_6$)} : Disponibilité temporelle de Florent (échelle 0-5, à maximiser).
    \item \textbf{Disponibilité Victor ($g_7$)} : Disponibilité temporelle de Victor (échelle 0-5, à maximiser).
\end{enumerate}

---

\section{Données collectées et normalisation}
\label{sec:donnees}

\subsection{Évaluations des alternatives (Matrice de performance)}
Le tableau~\ref{tab:perf_original} présente la matrice d'évaluation des 6 alternatives sur les 7 critères.

\begin{table}[H]
  \centering
  \caption{Tableau des performances initiales (matrice d'évaluation)}
  \label{tab:perf_original}
  \begin{tabular}{lccccccc}
    \toprule
    Alternative & $g_1$ (Coût) & $g_2$ (Jules) & $g_3$ (Florent) & $g_4$ (Victor) & $g_5$ (Disp. J) & $g_6$ (Disp. F) & $g_7$ (Disp. V) \\
    \midrule
    a1 (Me) & 75~€ & 4.0/10 & 9.0/10 & 5.0/10 & 5.0/5 & 5.0/5 & 5.0/5 \\
    a2 (Je) & 75~€ & 3.0/10 & 4.0/10 & 9.0/10 & 5.0/5 & 5.0/5 & 5.0/5 \\
    a3 (Ve) & 75~€ & 8.0/10 & 5.0/10 & 4.0/10 & 5.0/5 & 1.0/5 & 5.0/5 \\
    a4 (Sa) & 75~€ & 5.0/10 & 6.0/10 & 6.0/10 & 5.0/5 & 5.0/5 & 5.0/5 \\
    a5 (Di) & 75~€ & 9.0/10 & 3.0/10 & 7.0/10 & 5.0/5 & 5.0/5 & 1.0/5 \\
    a6 (3 Jours) & 185~€ & 7.3/10 & 4.7/10 & 5.7/10 & 5.0/5 & 2.0/5 & 2.0/5 \\
    \bottomrule
  \end{tabular}
\end{table}

\textbf{Justification des données :}
\begin{itemize}
    \item \textbf{Satisfaction artistique} : Établie selon l'intérêt des têtes d'affiche par jour. Le mercredi plaît particulièrement à Florent (9.0). Le jeudi plaît beaucoup à Victor (9.0). Le vendredi et le dimanche plaisent à Jules (respectivement 8.0 et 9.0). Le samedi est une journée de compromis neutre pour le groupe (Jules 5.0, Florent 6.0, Victor 6.0). La note du Pass 3 jours correspond à la moyenne des journées incluses (Vendredi, Samedi, Dimanche).
    \item \textbf{Disponibilité} : Jules est totalement libre (5/5 partout). Florent travaille le vendredi, sa disponibilité ce jour-là est très faible (1/5) et moyenne pour le pass (2/5). Victor ayant un repas de famille obligatoire le dimanche, sa disponibilité est de 1/5 ce jour-là et 2/5 pour le pass.
\end{itemize}

\subsection{Normalisation des critères}
Pour pouvoir appliquer des méthodes d'agrégation, il est nécessaire de ramener l'ensemble des critères sur une échelle de valeur commune $[0, 10]$ :
\begin{itemize}
    \item \textbf{Critère de Coût ($g_1$, à minimiser)} : Nous utilisons la transformation linéaire inverse :
    $$u_1(a) = 10 \times \frac{g_1^{\max} - g_1(a)}{g_1^{\max} - g_1^{\min}} = 10 \times \frac{185 - g_1(a)}{185 - 75}$$
    Un coût de 75~€ (le plus bas) donne une utilité de 10, et un coût de 185~€ (le plus élevé) donne une utilité de 0.
    \item \textbf{Critères artistiques ($g_2, g_3, g_4$, à maximiser)} : Déjà évalués sur $[0, 10]$, donc $u_j(a) = g_j(a)$ pour $j \in \{2, 3, 4\}$.
    \item \textbf{Critères de disponibilité ($g_5, g_6, g_7$, à maximiser)} : Évalués sur $[0, 5]$. On multiplie par 2 pour obtenir une note sur $[0, 10]$ : $u_j(a) = 2 \times g_j(a)$ pour $j \in \{5, 6, 7\}$.
\end{itemize}

Le tableau~\ref{tab:perf_normalized} présente la matrice des utilités normalisées.

\begin{table}[H]
  \centering
  \caption{Tableau des utilités normalisées (sur une échelle commune de 0 à 10)}
  \label{tab:perf_normalized}
  \begin{tabular}{lccccccc}
    \toprule
    Alternative & $u_1$ (Coût) & $u_2$ (Jules) & $u_3$ (Florent) & $u_4$ (Victor) & $u_5$ (Disp. J) & $u_6$ (Disp. F) & $u_7$ (Disp. V) \\
    \midrule
    a1 (Me) & 10.00 & 4.00 & 9.00 & 5.00 & 10.00 & 10.00 & 10.00 \\
    a2 (Je) & 10.00 & 3.00 & 4.00 & 9.00 & 10.00 & 10.00 & 10.00 \\
    a3 (Ve) & 10.00 & 8.00 & 5.00 & 4.00 & 10.00 & 2.00 & 10.00 \\
    a4 (Sa) & 10.00 & 5.00 & 6.00 & 6.00 & 10.00 & 10.00 & 10.00 \\
    a5 (Di) & 10.00 & 9.00 & 3.00 & 7.00 & 10.00 & 10.00 & 2.00 \\
    a6 (3 Jours) & 0.00 & 7.30 & 4.70 & 5.70 & 10.00 & 4.00 & 4.00 \\
    \bottomrule
  \end{tabular}
\end{table}

\subsection{Détermination des poids des critères}
La distribution de poids retenue est (la somme des poids est égale à $1.0$) :
\begin{itemize}
    \item \textbf{Coût ($g_1$)} : $w_1 = 0.25$.
    \item \textbf{Programmation artistique ($g_2, g_3, g_4$)} : $w_2 = w_3 = w_4 = 0.15$ chacun (total de $0.45$).
    \item \textbf{Disponibilité ($g_5, g_6, g_7$)} : $w_5 = w_6 = w_7 = 0.10$ chacun (total de $0.30$).
\end{itemize}

---

\section{Résolution par la méthode de la Somme Pondérée}
\label{sec:somme_ponderee}

La méthode de la Somme Pondérée calcule le score global de chaque alternative $a$ par la formule :
$$S(a) = \sum_{k=1}^{7} w_k \times u_k(a)$$

Le tableau~\ref{tab:ranking_ws} présente les scores globaux obtenus pour chaque alternative ainsi que le classement qui en découle.

\begin{table}[H]
  \centering
  \caption{Résultats et classement par la méthode de la Somme Pondérée}
  \label{tab:ranking_ws}
  \begin{tabular}{ccc}
    \toprule
    Rang & Alternative & Score global \\
    \midrule
    1 & a1 (Me) & 8.200 \\
    2 & a4 (Sa) & 8.050 \\
    3 & a2 (Je) & 7.900 \\
    4 & a5 (Di) & 7.550 \\
    5 & a3 (Ve) & 7.250 \\
    6 & a6 (3 Jours) & 4.455 \\
    \bottomrule
  \end{tabular}
\end{table}

\textbf{Analyse du classement :}
La somme pondérée place le Mercredi ($a_1$) en tête (score 8.200), suivi de près par le Samedi ($a_4$, 8.050) et le Jeudi ($a_2$, 7.900). Le Pass 3 jours ($a_6$) ferme la marche avec un faible score dû à son prix élevé et à l'effet cumulé des contraintes logistiques.

\textbf{Écueil de la méthode} : Elle est \textbf{compensatoire}. Le dimanche ($a_5$) et le vendredi ($a_3$) se classent devant le Pass 3 jours malgré les indisponibilités de Victor (dimanche) et Florent (vendredi). La méthode masque le fait que l'un des amis ne peut pas assister à ces journées.

---

\section{Résolution par la méthode de surclassement ELECTRE I}
\label{sec:electre}

La méthode ELECTRE I permet d'établir une relation de surclassement $aSb$ (« $a$ est au moins aussi bon que $b$ ») en combinant concordance et non-discordance (seuil de veto).

\subsection{Modélisation des seuils de veto}
Nous définissons des seuils de veto $v_k$ sur l'échelle $[0, 10]$ :
\begin{itemize}
    \item \textbf{Veto sur le Coût ($v_1 = 9.0$)}.
    \item \textbf{Veto sur l'intérêt artistique ($v_2 = v_3 = v_4 = 6.0$)}.
    \item \textbf{Veto sur la disponibilité ($v_5 = v_6 = v_7 = 5.0$)} : Si la disponibilité d'un membre baisse d'une valeur de 2.5 ou plus sur l'échelle 0-5 (soit 5.0 d'écart d'utilité normalisée), un veto bloque la comparaison.
\end{itemize}

\subsection{Indice de concordance $C(a_i, a_j)$}
L'indice de concordance $C(a_i, a_j)$ est défini par :
$$C(a_i, a_j) = \sum_{k: u_k(a_i) \ge u_k(a_j)} w_k$$

La matrice de concordance calculée est présentée dans le tableau~\ref{tab:concordance}.

\begin{table}[H]
  \centering
  \caption{Matrice de concordance $C(a_i, a_j)$}
  \label{tab:concordance}
  \begin{tabular}{ccccccc}
    \toprule
    & $a_1$ & $a_2$ & $a_3$ & $a_4$ & $a_5$ & $a_6$ \\
    \midrule
    $a_1$ & 1.00 & 0.85 & 0.85 & 0.70 & 0.70 & 0.70 \\
    $a_2$ & 0.70 & 1.00 & 0.70 & 0.70 & 0.85 & 0.70 \\
    $a_3$ & 0.60 & 0.75 & 1.00 & 0.60 & 0.60 & 0.75 \\
    $a_4$ & 0.85 & 0.85 & 0.85 & 1.00 & 0.70 & 0.85 \\
    $a_5$ & 0.75 & 0.60 & 0.75 & 0.75 & 1.00 & 0.75 \\
    $a_6$ & 0.40 & 0.40 & 0.35 & 0.25 & 0.35 & 1.00 \\
    \bottomrule
  \end{tabular}
\end{table}

\subsection{Indice de discordance et effet de veto}
Un veto est activé pour l'affirmation $a_i S a_j$ si $u_k(a_j) - u_k(a_i) > v_k$ pour au moins un critère $k$. Le tableau~\ref{tab:discordance} présente la matrice de discordance (1 = Veto activé, 0 = Pas de veto).

\begin{table}[H]
  \centering
  \caption{Matrice de discordance (1 = Veto activé, 0 = Pas de veto)}
  \label{tab:discordance}
  \begin{tabular}{ccccccc}
    \toprule
    & $a_1$ & $a_2$ & $a_3$ & $a_4$ & $a_5$ & $a_6$ \\
    \midrule
    $a_1$ & 0 & 0 & 0 & 0 & 0 & 0 \\
    $a_2$ & 0 & 0 & 0 & 0 & 0 & 0 \\
    $a_3$ & 1 & 1 & 0 & 1 & 1 & 0 \\
    $a_4$ & 0 & 0 & 0 & 0 & 0 & 0 \\
    $a_5$ & 1 & 1 & 1 & 1 & 0 & 0 \\
    $a_6$ & 1 & 1 & 1 & 1 & 1 & 0 \\
    \bottomrule
  \end{tabular}
\end{table}

L'analyse de cette matrice montre que :
\begin{itemize}
    \item $a_3$ (Vendredi) subit un veto à cause de l'indisponibilité de Florent le vendredi ($u_6 = 2$).
    \item $a_5$ (Dimanche) subit un veto à cause de l'indisponibilité de Victor le dimanche ($u_7 = 2$).
    \item Le Pass 3 jours ($a_6$) subit des vetoes systématiques à cause de son prix et de son incompatibilité avec Florent (le vendredi) et Victor (le dimanche).
    \item $a_1$ (Me), $a_2$ (Je) et $a_4$ (Sa) ne subissent **aucun veto** de la part des autres alternatives.
\end{itemize}

\subsection{Graphe de surclassement et noyaux multiples}
La relation de surclassement globale $a_i S a_j$ est définie par le franchissement d'un seuil de concordance $\gamma$ et l'absence de veto :
$$a_i S a_j \iff C(a_i, a_j) \ge \gamma \quad \text{et} \quad \text{Discordance}(a_i, a_j) = 0$$

Pour le seuil de concordance $\gamma = 0.60$, les alternatives $a_1$ (Me), $a_2$ (Je) et $a_4$ (Sa) surclassent toutes les alternatives du graphe et se surclassent mutuellement en formant un cycle de surclassement complet. La figure~\ref{fig:graphe_surclassement} illustre ce graphe de surclassement, qui se révèle particulièrement clair et structuré.

\begin{figure}[H]
  \centering
  \begin{tikzpicture}[
    node distance=2.5cm,
    main node/.style={circle, draw, minimum size=1.4cm, font=\bfseries}
  ]
    % Les trois alternatives dominantes (triangle supérieur)
    \node[main node] (a1) at (0, 3) {$a_1$};
    \node[main node] (a2) at (5, 3) {$a_2$};
    \node[main node] (a4) at (2.5, 1) {$a_4$};

    % Les trois alternatives dominées/veto (niveau inférieur)
    \node[main node] (a3) at (0, -1.5) {$a_3$};
    \node[main node] (a5) at (5, -1.5) {$a_5$};
    \node[main node] (a6) at (2.5, -4) {$a_6$};

    % Cycles symétriques (double-flèches) entre les gagnants a1, a2, a4
    \path[<->, >=Stealth, thick]
      (a1) edge (a2)
      (a1) edge (a4)
      (a2) edge (a4);
      
    % Flèches de surclassement vers les dominés (en gris pour la lisibilité)
    \path[->, >=Stealth, thick, gray]
      (a1) edge (a3)
      (a1) edge [bend right=15] (a6)
      (a2) edge (a5)
      (a2) edge [bend left=15] (a6)
      (a4) edge (a3)
      (a4) edge (a5)
      (a4) edge (a6);
      
    % Flèches complémentaires entre dominés
    \path[->, >=Stealth, thick, gray]
      (a3) edge (a6)
      (a5) edge (a6);
      
    % Flèches croisées de surclassement
    \path[->, >=Stealth, thick, gray]
      (a1) edge (a5)
      (a2) edge (a3);

  \end{tikzpicture}
  \caption{Graphe de surclassement ELECTRE I pour $\gamma = 0.60$}
  \label{fig:graphe_surclassement}
\end{figure}

Le calcul des noyaux du graphe donne trois noyaux distincts de taille minimale égale à 1 :
$$K_1 = \{a_1 \text{ (Mercredi)}\}, \quad K_2 = \{a_2 \text{ (Jeudi)}\}, \quad K_3 = \{a_4 \text{ (Samedi)}\}$$

Chacune de ces trois alternatives surclasse à elle seule toutes les autres alternatives hors de son noyau.

---

\section{Analyse comparative et discussion critique}
\label{sec:comparaison}

\subsection{L'effet de la compensation}
La Somme Pondérée masque les contraintes réelles sous l'effet de la compensation. Les alternatives $a_3$ (Vendredi) et $a_5$ (Dimanche) sont proposées en milieu de classement alors que Florent et Victor ne peuvent pas y assister. 
Dans ELECTRE I, les vetoes sur les disponibilités neutralisent immédiatement les options à contraintes fortes (Vendredi, Dimanche, Pass 3 jours), ce qui garantit le respect des contraintes absolues de chaque participant.

\subsection{Négociation collective et sens des noyaux multiples}
L'identification de trois noyaux uniques pour ELECTRE I ($\gamma = 0.60$) traduit les trois arbitrages fondamentaux de la décision de groupe :
\begin{enumerate}
    \item \textbf{Le choix orienté Florent ($a_1$ / Mercredi)} : Cette journée satisfait au maximum Florent (note artistique 9.0) et est tout à fait acceptable pour Victor et Jules, avec un coût de 75~€ et une pleine disponibilité.
    \item \textbf{Le choix orienté Victor ($a_2$ / Jeudi)} : Privilégie la journée favorite de Victor (note 9.0) tout en respectant le budget et l'emploi du temps du groupe.
    \item \textbf{Le choix du consensus équitable ($a_4$ / Samedi)} : Journée moyenne pour tout le monde (Jules 5.0, Florent 6.0, Victor 6.0) mais sans déception majeure et sans contraintes.
\end{enumerate}

Le choix final au sein du groupe peut ainsi s'opérer par une négociation autour de ces trois options dominantes. Par exemple, si Jules (dont les jours préférés Vendredi et Dimanche sont exclus par veto) doit trancher, il aura intérêt à préférer le samedi ($a_4$, note 5.0) comme compromis équitable ou le mercredi ($a_1$, note 4.0) par solidarité avec Florent.

---

\section{Évaluation critique de l'usage de l'IA Générative}
\label{sec:evaluation_ia}

\subsection{Nature et volume des requêtes}
L'IA a été sollicitée pour aider à la modélisation mathématique du problème, au développement du script Python pur et à la structuration du fichier LaTeX, soit 4 requêtes principales au total.

\subsection{Retours critiques et corrections apportées}
Nous avons apporté une modification importante au problème initial. Pour éviter un graphe de surclassement trop encombré (avec 9 alternatives, le graphe TikZ devenait illisible et surchargeait le rapport), nous avons simplifié la modélisation à 6 alternatives clés (les 5 jours individuels et le Pass 3 jours). Cette simplification a permis d'obtenir une figure~\ref{fig:graphe_surclassement} extrêmement lisible et élégante tout en conservant intacte la richesse de la confrontation des méthodes (compensation vs veto et noyaux multiples). Nous avons également dû corriger le code Python généré qui utilisait initialement \texttt{numpy}.

\subsection{Avantages et inconvénients}
L'IA s'est révélée précieuse pour la mise en forme du code LaTeX (TikZ et tableaux professionnels) et le calcul automatisé des matrices, nous évitant les erreurs manuelles de calcul de concordance. Cependant, elle requiert une grande vigilance théorique pour corriger ses analyses génériques et s'assurer du respect des spécificités du cours.

\end{document}
